\section{傅里叶变换的部分相关概念及性质}

\subsection{傅里叶级数}
利用一套完备的正交函数系，我们可以将一个周期函数展开成该正交函数系的线性组合。对于周期函数来说，最常用的完备正交函数系就是三角函数系$\{\cos(n\omega x),\sin(n\omega x)\}$。下面我们给出傅里叶级数定理：
\begin{theorem}[傅里叶级数定理]
当一个周期函数$f(x)$满足迪利克雷（Dirichlet）条件,即在区间${[-\frac{T}{2},\frac{T}{2}]}$上满足
\begin{itemize}
    \item 1. f(x)在区间${[-\frac{T}{2},\frac{T}{2}]}$上连续或者只有有限个第一间断点
    \item 2. f(x)在区间${[-\frac{T}{2},\frac{T}{2}]}$上只有有限个极值点
\end{itemize}
$f(x)$就可以展开成傅里叶级数
\[f(x)=\frac{a_0}{2}+\sum_{n=1}^{+\infty}(a_n\cos(n\omega x)+b_n\sin(n\omega x))\]
这里我们利用周期$T$定义频率$\omega$
\[\omega=\frac{2\pi}{T}\]
则傅立叶展开的系数可以表示成
\[a_0=\frac{2}{T}\int_{-\frac{T}{2}}^{\frac{T}{2}}f(x)\dd{x}, \quad a_n=\frac{2}{T}\int_{-\frac{T}{2}}^{\frac{T}{2}}f(x)\cos(n\omega x)\dd{x}, \quad b_n=\frac{2}{T}\int_{-\frac{T}{2}}^{\frac{T}{2}}f(x)\sin(n\omega x)\dd{x}\]
\end{theorem}
将$\cos(n\omega x)$、$\sin(n\omega x)$写成复指数形式，则傅里叶级数可以写成以下形式
\[f(x)=\frac{a_0}{2}+\sum_{n=1}^{+\infty}(\frac{a_n-ib_n}{2}e^{in\omega x}+\frac{a_n+ib_n}{2})e^{-in\omega x}\]
定义一套新的系数$\{c_n\}$：
\[c_0=\frac{a_0}{2}=\frac{1}{T}\int_{-\frac{T}{2}}^{\frac{T}{2}}f(x)\dd{x}\]
\[c_n=\frac{a_n-ib_n}{2}=\frac{1}{T}(\int_{-\frac{T}{2}}^{\frac{T}{2}}f(x)(\cos(n\omega x)\dd{x}-i\sin(n\omega x))\dd{x})=\frac{1}{T}\int_{-\frac{T}{2}}^{\frac{T}{2}}f(x)e^{-in\omega x}\dd{x}\]
\[c_{-n}=\frac{a_n+ib_n}{2}=\frac{1}{T}(\int_{-\frac{T}{2}}^{\frac{T}{2}}f(x)(\cos(n\omega x)\dd{x}+i\sin(n\omega x))\dd{x})=\frac{1}{T}\int_{-\frac{T}{2}}^{\frac{T}{2}}f(x)e^{in\omega x}\dd{x}\]
即可得到复指数形式的傅立叶级数定理：
\begin{theorem}[复指数形式的傅里叶级数定理]
\[f(x)=c_0+\sum_{n=1}^{+\infty}(c_ne^{in\omega x}+c_{-n}e^{-in\omega x})=\sum_{n=-\infty}^{+\infty}c_ne^{in\omega x}, \quad c_n=\frac{1}{T}\int_{-\frac{T}{2}}^{\frac{T}{2}}f(x)e^{-in\omega x}\dd{x}\]
或者简单写成
\[f(x)=\sum_{n=-\infty}^{+\infty}(\frac{1}{T}\int_{-\frac{T}{2}}^{\frac{T}{2}}f(y)e^{-i\omega_n y}\dd{y})e^{i\omega_n x}\]
\end{theorem}
对于周期函数，其傅立叶变换会退化会离散的傅立叶级数。

\newpage
\subsection{傅里叶积分定理与傅里叶变换}
在实际的应用中，我们得到的信号曲线不一定是一个周期函数，那么我们有没有办法将一个非周期的函数展开成傅里叶级数呢？对于一个定义在实数轴上的非周期的函数，我们可以认为它的周期是无穷大，对此我们将把能展开成傅里叶级数的周期函数延拓到整个实数轴上，再将其周期取为无穷大。

将$f(x)$按周期$T$在实数域上做延拓的到函数$f_T(x)$：
\[ f_T(x) := \left\{
\begin{array}{rl}
f(x) & \text{if } x \in [-\frac{T}{2},\frac{T}{2}]\\
f_T(x \pm T) & \text{if } x \notin [-\frac{T}{2},\frac{T}{2}] 
\end{array} \right.\]
把周期扩大到无穷$T \rightarrow \infty$时，我们可以得到：
\[\lim_{T \rightarrow \infty}{f_T(x)=f(x)}\]
对于一个可以展开成傅里叶级数且其周期$T$无穷大的函数，其傅里叶级数可以写成
\[f(x)=\lim_{T \rightarrow \infty}\frac{1}{T}\sum_{n=-\infty}^{+\infty}(\int_{-\frac{T}{2}}^{\frac{T}{2}}f(y)e^{-in\omega y}\dd{y})e^{in\omega x}=\lim_{\omega \rightarrow 0}\frac{\omega}{2\pi}\sum_{n=-\infty}^{+\infty}(\int_{-\frac{T}{2}}^{\frac{T}{2}}f(y)e^{-in\omega y}\dd{y})e^{in\omega x}\]
由积分的定义式上式最终可以得到
\[f(x)=\frac{1}{2\pi}\int_{-\infty}^{+\infty}(\int_{-\infty}^{+\infty}f(y)e^{-i\omega y}\dd{y})e^{i\omega x}\dd{\omega}\]
\begin{theorem}[傅里叶积分公式]
当$f(x)$在${(-\infty,+\infty)}$满足：
\begin{itemize}
    \item 1. 在任一有限区间上满足迪利克雷条件；
    \item 2. 在${(-\infty,+\infty)}$上绝对可积（即$\int_{-\infty}^{+\infty}|f(x)|\dd{x}$收敛）；
\end{itemize}
当$f(x)$在点$x=x_0$处连续时
\[\frac{1}{2\pi}\int_{-\infty}^{+\infty}(\int_{-\infty}^{+\infty}f(y)e^{-i\omega y}\dd{y})e^{i\omega x_0}\dd{\omega}=f(x_0)\]
当$f(x)$在点$x=x_0$处不连续时
\[\frac{1}{2\pi}\int_{-\infty}^{+\infty}(\int_{-\infty}^{+\infty}f(y)e^{-i\omega y}\dd{y})e^{i\omega x_0}\dd{\omega}=\frac{1}{2}\left ( \lim_{x \rightarrow x_0^-}{f(x)}+\lim_{x \rightarrow x_0^+}{f(x)} \right )\]
\end{theorem}
\begin{definition}[傅里叶变换与傅里叶逆变换]
设$f(t)$在${(-\infty,+\infty)}$上满足傅里叶积分定理的条件，可定义以下积分变换
\[F(\omega)=\int_{-\infty}^{+\infty}f(t)e^{-i\omega t}\dd{t}, \quad f(t)=\frac{1}{2\pi}\int_{-\infty}^{+\infty}F(\omega)e^{i\omega t}\dd{\omega}\]
称$F(\omega)$为$f(t)$的象函数，称$f(t)$到$F(\omega)$的积分变换为傅里叶变换，,$f(t)$为$F(\omega)$的象原函数，称$F(\omega)$到$f(t)$的积分变换为傅里叶逆变换式，记为$F(\omega)=\mathcal{F}[f(t)], \ f(t)=\mathcal{F}^{-1}[F(\omega)]$。
\end{definition}
傅立叶变换中一个重要的应用对象是$\delta$函数（单位脉冲函数），是英国物理学家狄拉克(Dirac)在20世纪20年代引人的，用于描述瞬间或空间几何点上的物理量。例如，瞬时的冲击力、脉冲电流或电压等急速变化的物理量，以及质点的质量分布、点电荷的电量分布等在空间或时间上高度集中的物理量。
\begin{definition}[单位脉冲函数]
设函数$\delta(t)$为一维空间中自变量为时间t的函数，若它满足:
\[\delta(t)=\left\{
\begin{array}{rl}0 & \text{if } t \neq 0\\
1 & \text{if } t=0
\end{array} \right., \quad \int_{-\infty}^{+\infty}\delta(t)\dd{t}=1\]
则称$\delta(t)$为单位脉冲函数，或简称为$\delta$函数。
\end{definition}
$\delta$函数有如下重要的性质（\textbf{筛选性质}）：
\[\int_{-\infty}^{+\infty}\delta(t-t_0)f(t)\dd{t}=f(t_0)\]
利用这个性质，我们可以对单位脉冲函数求其傅里叶变换及傅里叶逆变换：
\[F(\omega)=\mathcal{F}[\delta(t-t_0)]=\int_{-\infty}^{+\infty}\delta(t-t_0)e^{-i\omega t}\dd{t}=e^{-i\omega t_0}\]
令$F(\omega)=2\pi \delta(\omega-\omega_0)$则有： 
\[f(t)=\mathcal{F}^{-1}[2\pi \delta(\omega-\omega_0)]=\frac{1}{2\pi}\int_{-\infty}^{+\infty}2\pi \delta(\omega-\omega_0)e^{i\omega t}\dd{\omega}=e^{i\omega_0 t}\]
这是我们得到两对傅里叶变换对：$\delta(t-t_0)$和$e^{-i\omega t_0}$，$e^{i\omega_0 t}$和$2\pi \delta(\omega-\omega_0)$。可以看出$\delta$函数与平面波是傅立叶变换对。特别的，我们令$t_0=0,\omega_0=0$，则有$\delta(t)$和$1$，$1$和$2\pi \delta(\omega)$。
稍作推广我们可以得到
\[\mathcal{F}[\cos(\omega_0t)]=\pi[\delta(\omega+\omega_0)+\delta(\omega-\omega_0)], \quad \mathcal{F}[\sin(\omega_0t)]=i\pi[\delta(\omega+\omega_0)-\delta(\omega-\omega_0)]\]
\begin{proposition}[傅里叶变换的基本性质]
设$a,b \in \mathbb{C}$，$f,f_1,f_2$是满足傅立叶积分定理的复值函数，则傅立叶变换具有以下性质:
\begin{itemize}
\item 1. 线性
\[\mathcal{F}[af_1(t)+bf_2(t)]=a\mathcal{F}[f_1(t)]+b\mathcal{F}[f_2(t)], \quad \mathcal{F}^{-1}[af_1(t)+bf_2(t)]=a\mathcal{F}^{-1}[f_1(t)]+b\mathcal{F}^{-1}[f_2(t)]\]
\item 2. 位移性质
\[\mathcal{F}[f(t \pm t_0)]=e^{\pm i\omega t_0}\mathcal{F}[f(t)], \quad \mathcal{F}^{-1}[f(t \pm t_0)]=e^{\mp i\omega t_0}\mathcal{F}^{-1}[f(t)]\]
\item 3. 微分性质：若$f \in C^n(-\infty,+\infty)$（可以推广到有限个间断点），且$\lim_{|t|\rightarrow +\infty}f^{(k)}(t)=0$:
\[\mathcal{F}[f^{(n)}(t)]=(i\omega)^n \mathcal{F}[f(t)], \quad \frac{\mathrm {d}^n}{\mathrm {d} \omega^n} \mathcal{F}[f(t)]=(-i)^n \mathcal{F}[t^nf(t)]\]
\item 4. 积分性质：当$t \rightarrow +\infty$时，$\int_{-\infty}^{t}f(t)\dd{t}<+\infty$，则
\[ \mathcal{F}[\int_{-\infty}^{t}f(t)\dd{t}]=\frac{1}{i\omega}\mathcal{F}[f(t)]+\pi\delta(\omega)\int_{-\infty}^{\infty}f(t)\dd{t}\]
\item 5. 乘积定理：若$F_1(\omega)=\mathcal{F}[f_1(t)]$，$F_2(\omega)=\mathcal{F}[f_2(t)]$，则
\[\int_{-\infty}^{+\infty}f_1^{\dagger}(t)f_2(t)\dd{t}=\frac{1}{2\pi}\int_{-\infty}^{+\infty}F_1^{\dagger}(\omega)F_2(\omega)\dd{\omega}, \quad \int_{-\infty}^{+\infty}f_1(t)f_2^{\dagger}(t)\dd{t}=\frac{1}{2\pi}\int_{-\infty}^{+\infty}F_1(\omega)F_2^{\dagger}(\omega)\dd{\omega}\]
\item 6.能量积分(Parseval等式)
\[\int_{-\infty}^{+\infty}[f(t)]^2\dd{t}=\frac{1}{2\pi}\int_{-\infty}^{+\infty}[F(\omega)]^2\dd{\omega}:=\frac{1}{2\pi}\int_{-\infty}^{+\infty}S(\omega)\dd{\omega}\]
\end{itemize}
\end{proposition}
\subsection{卷积与相关函数}
我们可以定义卷积运算及相关函数
\begin{definition}[卷积]
若已知函数$f_1(t),f_2(t)$，则定义卷积为
\[f_1(t)\ast f_2(t):=\int_{-\infty}^{+\infty}f_1(\tau)f_2(t-\tau)\dd{\tau}\] 
\end{definition}
\begin{definition}[相关函数]
给定两函数$f_1(t)$、$f_2(t)$,定义相关函数（$f_1(t)=f_2(t)$时，称$R(\tau)$为自相关函数）：
\[R(\tau):=\int_{-\infty}^{+\infty}f_1(t)f_2(t+\tau)\dd{t}\]
\end{definition}
\begin{proposition}[卷积的性质]
\begin{itemize}
\item 1. 交换律：$f_1(t)\ast f_2(t)=f_2(t)\ast f_1(t)$
\item 2. 结合律：$(f_1(t) \ast f_2(t)) \ast f_3(t)=f_1(t) \ast (f_2(t) \ast f_3(t))$
\item 3. 线性性：$f_1(t) \ast (af_2(t)+bf_3(t))=af_1(t) \ast f_2(t)+f_1(t) \ast bf_3(t),  \ a,b\in\mathbb{C}$
\item 4. 微分：
\[\frac{\mathrm {d}}{\mathrm {d}x}[f_1(t)\ast f_2(t)]=[\frac{\mathrm {d}}{\mathrm {d}x}f_1(t)]\ast f_2(t)=f_1(t)\ast\frac{\mathrm {d}}{\mathrm {d}x}[ f_2(t)]\]
\item 5. 卷积不等式：
\[|f_1(t)\ast f_2(t)| \leqslant |f_1(t)|\ast |f_2(t)|\]
\item 6. $\delta$函数的卷积性质：
\[f(t) \ast \delta(t)=f(t)\]
\end{itemize}
\end{proposition}
卷积最重要的性质是卷积定理
\begin{theorem}[卷积定理]
    若$f_1(t)$和$f_2(t)$满足傅里叶积分定理的条件，则有
\[\mathcal{F}[f_1(t) \ast f_2(t)]=\mathcal{F}[f_1(t)] \cdot \mathcal{F}[f_2(t)]\]
\end{theorem}

\section{傅里叶变换的应用}
% 2 FFT算法的原理
\subsection{Borwein积分}
下面按照 \href{https://www.bilibili.com/video/BV18e4y1u7BH/}{3Blue1Brown 关于 Borwein 积分的视频} 的思路来组织这一节。先看一列积分：
\[I_n=\int_0^{+\infty}\prod_{k=0}^{n}\frac{\sin\left(x/(2k+1)\right)}{x/(2k+1)}\dd{x}\]
数值上我们会发现
\[I_0=I_1=\cdots=I_6=\frac{\pi}{2}, \qquad I_7=\frac{\pi}{2}-2.31\times 10^{-11}\]
也就是说，把$\sin(x/3)/(x/3),\sin(x/5)/(x/5),\ldots,\sin(x/13)/(x/13)$依次乘进去以后，积分值始终不变；但再乘上$\sin(x/15)/(x/15)$之后，它才第一次偏离$\pi/2$，而且偏离量还极其微小。更夸张的是，再考虑
\[\int_0^{+\infty}2\cos(x)\frac{\sin(x)}{x}\frac{\sin(x/3)}{x/3}\cdots\frac{\sin(x/111)}{x/111}\dd{x}=\frac{\pi}{2}\]
但如果在这个乘积后面再乘上$\sin(x/113)/(x/113)$，结果才会第一次小于$\pi/2$，而偏离量大约只有$2.3324\times 10^{-138}$。难怪最早做数值实验的人会怀疑计算机是不是出了bug。

如果只盯着这些积分本身，很难一眼看出为什么会在$15$或$113$这样的数字上突然失效。视频中的核心思路是：先把这个问题翻译成一个更容易看清几何结构的问题，也就是\textbf{矩形平台经过反复滑动平均后，平台什么时候会被“吃掉”}。

先定义矩形函数
\[\operatorname{rect}(x):=\left\{
\begin{array}{rl}
1 & \text{if } |x|<\frac{1}{2}\\
\frac{1}{2} & \text{if } |x|=\frac{1}{2}\\
0 & \text{if } |x|>\frac{1}{2}
\end{array}\right.\]
它在中间有一个长度为$1$的平台。再定义一列函数：
\[f_0(x)=\operatorname{rect}(x), \qquad f_{m+1}(x)=\frac{1}{a_{m+1}}\int_{x-a_{m+1}/2}^{x+a_{m+1}/2}f_m(y)\dd y\]
其中
\[a_m=\frac{1}{2m+1}\]
于是$f_1$是把$f_0$在宽度为$1/3$的窗口上做滑动平均，$f_2$是再用宽度为$1/5$的窗口做一次滑动平均，之后依次用$1/7,1/9,\ldots$的窗口继续平均下去。

这个模型最关键的观察是：只要滑动窗口还能够完全放进上一步函数的平台里，那么新函数在原点附近就仍然会保留一个高度为$1$的平台。因此平台长度会依次减少
\[1,\qquad 1-\frac{1}{3},\qquad 1-\frac{1}{3}-\frac{1}{5},\qquad 1-\frac{1}{3}-\frac{1}{5}-\frac{1}{7},\ \ldots\]
也就是说，在
\[\frac{1}{3}+\frac{1}{5}+\cdots+\frac{1}{2n+1}<1\]
之前，原点$x=0$始终还落在平台内部，因此$f_n(0)=1$。但一旦这个和第一次超过$1$，平台就在原点附近消失，$f_n(0)$便会第一次小于$1$。直接计算可得
\[\frac{1}{3}+\frac{1}{5}+\frac{1}{7}+\frac{1}{9}+\frac{1}{11}+\frac{1}{13}<1,\qquad
\frac{1}{3}+\frac{1}{5}+\frac{1}{7}+\frac{1}{9}+\frac{1}{11}+\frac{1}{13}+\frac{1}{15}>1\]
因此第一次失效恰好发生在$15$。而且因为平台消失前，平台边缘附近的函数值早就已经被多次平均得非常接近$1$，所以即便失效，$f_n(0)$也只会比$1$小极其微弱的一点点。这正解释了为什么$I_7$只是比$\pi/2$略微小了$10^{-11}$量级。

对于带$2\cos x$的版本，同样的阈值为什么会从$15$推迟到$113$呢？直观地说，视频里的答案是：\textbf{起始平台变长了}。这个结论可以从恒等式
\[2\cos(\pi t)\operatorname{sinc}_{\pi}(t)=2\operatorname{sinc}_{\pi}(2t),\qquad \operatorname{sinc}_{\pi}(t):=\frac{\sin(\pi t)}{\pi t}\]
看出来。它说明在一个更自然的归一化变量里，原来“$2\cos x$乘上第一个sinc因子”的效果，相当于把最初的矩形平台长度从$1$扩大到$2$。于是平台长度的判据从“倒数和是否超过$1$”变成了“倒数和是否超过$2$”：
\[\frac{1}{3}+\frac{1}{5}+\cdots+\frac{1}{111}<2,\qquad
\frac{1}{3}+\frac{1}{5}+\cdots+\frac{1}{113}>2\]
因此新的阈值就自然落在$113$，这和视频中的数值实验完全一致。

剩下的问题是：这个“滑动平均吃掉平台”的模型为什么会和原来的积分问题一模一样？这正是傅里叶变换与卷积定理发挥作用的地方。为了和视频中的图像解释保持一致，本小节暂时采用归一化傅里叶变换
\[\widehat{f}(\xi)=\int_{-\infty}^{+\infty}f(x)e^{-2\pi i x\xi}\dd x\]
在这个约定下有三条最关键的事实：
\[\widehat{\operatorname{sinc}_{\pi}(ax)}(\xi)=\frac{1}{a}\operatorname{rect}\left(\frac{\xi}{a}\right),\qquad
\int_{-\infty}^{+\infty}f(x)\dd x=\widehat{f}(0),\qquad
\widehat{fg}=\widehat{f}\ast\widehat{g}\]
其中$\ast$是通常意义下的卷积。记
\[r_a(\xi):=\frac{1}{a}\operatorname{rect}\left(\frac{\xi}{a}\right)\]
则$r_a$是一个宽度为$a$、总面积为$1$的矩形窗，而且
\[(f\ast r_a)(x)=\frac{1}{a}\int_{x-a/2}^{x+a/2}f(y)\dd y\]
恰好就是宽度为$a$的滑动平均。

现在对$I_n$做代换$x=\pi t$，可得
\[I_n=\frac{\pi}{2}\int_{-\infty}^{+\infty}\prod_{k=0}^{n}\operatorname{sinc}_{\pi}\!\left(\frac{t}{2k+1}\right)\dd t\]
于是
\[\begin{aligned}
I_n
&=\frac{\pi}{2}\left[\widehat{\prod_{k=0}^{n}\operatorname{sinc}_{\pi}\!\left(\frac{t}{2k+1}\right)}\right](0)\\
&=\frac{\pi}{2}\left(r_1\ast r_{1/3}\ast r_{1/5}\ast\cdots\ast r_{1/(2n+1)}\right)(0)
\end{aligned}\]
这就表明：\textbf{Borwein积分其实正是在问，一连串矩形窗做滑动平均以后，原点处的函数值是多少。} 因而$I_n$等于$\pi/2$还是略小于$\pi/2$，完全取决于那个平台在原点处是否还存在。至此，原本看似神秘的“前几项始终等于$\pi/2$，直到某一步才突然失效”的现象，就被翻译成了一个非常直观的几何过程。
